\section{Identifiability and cell-total consistency}
\label{sec:identifiability}

We port \cref{thm:bg-id} to ST and derive a cell-total consistency
corollary.

\subsection{The rank condition}

Let the spot-by-cell-type composition matrix $P \in [0,1]^{S \times K}$
have entries $P_{sk} = p_{s,k}$, with rows summing to $1$. For each
gene $j$, the spot-level mean expression is $\E[X_{sj}/N_s] = \sum_k
P_{sk} \mu_{j,k}$. In matrix form, $\bm{m}_j = P \bm{\mu}_j$ where
$\bm{m}_j \in \R^S$ is the vector of per-spot mean expressions and
$\bm{\mu}_j \in \R^K$ is the per-cell-type mean expression vector
for gene $j$.

\begin{theorem}[Identifiability of cell-type-specific expression]
\label{thm:rank}
Assume Poisson sampling $X_{sj} \sim
\mathrm{Poisson}\bigl(N_s \sum_k p_{s,k} \mu_{j,k}\bigr)$ with the
composition matrix $P$ known (e.g.\ supplied by an external
reference). The cell-type-specific mean expression vector
$\bm{\mu}_j \in \R_{\geq 0}^K$ is identifiable from spot-level data
$\{X_{sj}\}_{s=1\ldots S}$ if and only if $P$ has full column rank
$K$.
\end{theorem}

\begin{proof}
Identifiability of $\bm{\mu}_j$ from the joint distribution of
$\{X_{sj}\}_s$ is equivalent to identifiability of the spot-level
mean vector $\bm{m}_j = P\bm{\mu}_j$, because the Poisson family
is parametrized by its mean (so two parameter vectors yielding the
same mean vector yield the same joint distribution).

\emph{Necessary direction.} Suppose $P$ has column rank less than
$K$, so there exists $\bm{v} \neq \bm{0}$ with $P\bm{v} = \bm{0}$.
For any candidate parameter $\bm{\mu}_j$, the perturbed
$\bm{\mu}_j' = \bm{\mu}_j + \bm{v}$ (chosen so $\bm{\mu}_j'$ remains
in $\R_{\geq 0}^K$, possible whenever $\bm{\mu}_j$ has all positive
entries) satisfies $P\bm{\mu}_j' = P\bm{\mu}_j = \bm{m}_j$. Both
parameters induce identical spot-level distributions of $\{X_{sj}\}$,
so $\bm{\mu}_j$ is not identifiable.

\emph{Sufficient direction.} Suppose $P$ has full column rank $K$.
Then $P^\top P$ is invertible, and the map $\bm{\mu}_j \mapsto
\bm{m}_j = P\bm{\mu}_j$ is injective on $\R^K$: distinct parameters
yield distinct mean vectors, and identifiability of $\bm{\mu}_j$
follows. The closed-form left inverse
$\bm{\mu}_j = (P^\top P)^{-1} P^\top \bm{m}_j$ recovers the
parameter from the mean exactly.
\end{proof}

\begin{remark}[Extensions]
\label{rmk:rank-extensions}
Three extensions of \cref{thm:rank} are immediate. (i) If $N_s$
varies across spots, the proof is unchanged: identifiability is a
property of the parametric family, not the sample size; consistent
estimation of $\bm{m}_j$ requires $\min_s N_s$ or the harmonic mean
to grow, which is the standard Poisson-MLE regularity condition.
(ii) Bernoulli dropout, in which $X_{sj}$ is observed as zero with
some cell-type-independent probability $\pi$, is absorbed into the
mean: $\E[X_{sj}/N_s] = (1-\pi) \sum_k p_{s,k} \mu_{j,k}$, and the
rank condition on $P$ remains necessary and sufficient (up to a
known constant). (iii) Cell-type-dependent dropout $\pi_k$
violates condition C2 of the masking framework and is the regime
treated by \cite{towell2026mdrelax}; the rank condition is no
longer sufficient and identifiability degrades quantitatively in
the gap between the maximum and minimum $\pi_k$.
\end{remark}

\Cref{thm:rank} is the ST analogue of \cref{thm:bg-id}. It clarifies
the role of single-cell-resolution platforms: a single cell measured
at single-cell resolution contributes a row of $P$ that is a unit
vector $\bm{e}_k$ (composition $100\%$ type $k$), maximally
informative for column $k$ of $P$. The classical reliability statement
``singletons restore identifiability'' becomes ``single-cell-resolution
spots restore identifiability''.

\subsection{When does the rank condition fail?}

In purely sequencing-based ST (no MERFISH/seqFISH+ integration), $P$
typically has full column rank $K$ \emph{across spots} as long as
spots cover the full cellular diversity of the tissue. The condition
fails when:
\begin{itemize}
\item Two cell types always co-occur in the same proportion (e.g.,
  endothelial and pericyte cells in capillary regions), making their
  contributions inseparable in $P$.
\item The number of distinct cell types $K$ exceeds the number of
  distinct $\bm{p}_s$ vectors observed (rare in practice for
  whole-tissue assays).
\item Cell types are functionally identical with respect to the
  measured genes; in this case the framework is correctly reporting
  that no method can distinguish them.
\end{itemize}

\subsection{Cell-total consistency}

\begin{theorem}[Cell-total consistency]
\label{thm:cell-total-st}
Let $(\hat{\bm{\mu}}_j, \hat{P})$ be an MLE under the spot-level
Poisson likelihood with mean $N_s \sum_k p_{s,k}\mu_{j,k}$, in which
both $\bm{\mu}_j$ and $\bm{p}_s$ are unconstrained at the optimum
(no prior, no penalty, no spatial-smoothness regularization that
ties the two together). At any interior MLE, for each spot $s$ and
each gene $j$,
\begin{equation}
  \sum_k \hat{P}_{sk} \hat{\mu}_{j,k} = \bar{X}_{sj},
  \label{eq:cell-total-st}
\end{equation}
where $\bar{X}_{sj} = X_{sj}/N_s$ is the empirical mean. The
identity holds exactly (up to optimization tolerance); boundary
cases ($\hat P_{sk} = 0$ for some $k$, or $\hat\mu_{j,k} = 0$ for
some $(j,k)$) reduce to identities on the active support.
\end{theorem}

\begin{proof}[Sketch]
Parallels the ZINB cell-total identity in the scRNA-seq precursor
paper. Under the unconstrained MLE assumption, the score equation
with respect to $\hat\mu_{j,k}$ at an interior optimum reads
$\sum_s \hat P_{sk} (\bar X_{sj} - \sum_{k'} \hat P_{sk'}
\hat\mu_{j,k'}) = 0$ for each $(j,k)$, and the score equation with
respect to $\hat P_{sk}$ reads $\sum_j \hat\mu_{j,k} (\bar X_{sj} -
\sum_{k'} \hat P_{sk'} \hat\mu_{j,k'}) = 0$ for each $(s,k)$. Writing
the residual matrix $R_{sj} = \bar X_{sj} - \sum_{k'} \hat P_{sk'}
\hat\mu_{j,k'}$ and the full-transcriptome estimator matrix
$\hat U = [\hat\mu_{j,k}] \in \R^{J \times K}$ (distinct from the
marker-restricted $M \in \R^{|\mathcal{M}| \times K}$ of
\cref{sec:methodology}), both score equations imply $\hat P^\top R
= 0$ and $R \hat U = 0$ on the active support. When both $\hat P$
and $\hat U$ have full column rank on their respective active
supports (the conditions of \cref{thm:rank} applied to the joint
estimator), the only common-kernel solution is $R = 0$ on the
active support, giving \eqref{eq:cell-total-st} exactly at any
interior MLE.
\end{proof}

\textbf{Scope and practical implication.} \Cref{thm:cell-total-st}
applies to unregularized MLE-based deconvolution: this covers naive
RCTD and the unregularized limit of joint NMF. It does \emph{not}
cover Cell2location's Bayesian-hierarchical prior or CARD's
spatial-smoothness penalty, which couple $\bm{p}_s$ across spots and
break the unconstrained-optimum assumption. The implication for
unregularized methods is sharp: disagreement on per-cell-type
estimates is invisible in marginal-fit checks. For regularized
methods, the closer the regularization is to identity, the more
\eqref{eq:cell-total-st} continues to hold approximately. Diagnosing
per-cell-type bias requires single-cell-resolution data, marker-gene
calibration, or external biological priors.
